\begin{theorem}[无外加寄存器、置换自然性约束下的二值跳跃结构唯一性]\label{thm:bdry-binary-jump-orientation-functor-uniqueness}
令 \(k\ge 2\)，并令 \(\mathbf{FinSet}_k^{\cong}\) 表示由所有 \(k\) 元有限集合与其双射组成的群胚（任意两对象同构，且 \(\Aut([k])\cong \mathfrak S_k\)）。
令 \(\mathbf{Tors}(\ZZ_2)\) 表示 \(\ZZ_2\)-torsor（两点集合上的自由传递 \(\ZZ_2\) 作用）与其同构构成的群胚。
称一个\emph{二值跳跃结构}为一函子
\[
\mathcal{J}:\mathbf{FinSet}_k^{\cong}\longrightarrow \mathbf{Tors}(\ZZ_2).
\]
则当 \(k\ge 3\) 时，除平凡函子外，在自然同构意义下存在且仅存在一个非平凡 \(\mathcal{J}\)；
其可由\emph{取向函子}给出：对任意 \(k\) 元集合 \(S\)，定义
\[
\mathrm{Or}(S):=\mathrm{Bij}([k],S)\big/\mathfrak A_k,
\]
其中 \(\mathfrak A_k\subset \mathfrak S_k\) 为交错群，并以 \(\mathfrak S_k\) 在 \(\mathrm{Bij}([k],S)\) 上的预合成作用诱导 \(\mathrm{Or}(S)\) 的 \(\ZZ_2\)-torsor 结构；其作用恰经由符号同态 \(\mathrm{sgn}:\mathfrak S_k\to\ZZ_2\) 因子化。
\end{theorem}

\begin{proof}
取基对象 \(S_0=[k]\)。
任一 \(\ZZ_2\)-torsor 的自同构群自然同构于 \(\ZZ_2\)（由“平移”给出），因此对任意函子 \(\mathcal{J}\)，由 \(\Aut(S_0)=\mathfrak S_k\) 在 \(\mathcal{J}(S_0)\) 上的函子性作用得到群同态
\[
\chi_{\mathcal{J}}:\mathfrak S_k\to \Aut(\mathcal{J}(S_0))\cong \ZZ_2.
\]
反之，给定任意群同态 \(\chi:\mathfrak S_k\to\ZZ_2\)，取一固定两点 \(\ZZ_2\)-torsor \(T\) 并令 \(\mathfrak S_k\) 通过 \(\chi\) 作用于 \(T\)；
再利用 \(\mathbf{FinSet}_k^{\cong}\) 的连通性与函子性，可沿任意双射将该作用在所有对象上传输，从而得到一函子 \(\mathcal{J}\)，且其同构类仅依赖于 \(\chi\)。
由此 \(\mathcal{J}\) 的自然同构类与 \(\mathrm{Hom}(\mathfrak S_k,\ZZ_2)\) 一一对应。
\par
当 \(k\ge 3\) 时 \(\mathfrak S_k^{\mathrm{ab}}\cong \ZZ_2\)，故 \(\mathrm{Hom}(\mathfrak S_k,\ZZ_2)\) 仅含平凡同态与符号同态两者，从而非平凡 \(\mathcal{J}\) 在自然同构意义下唯一。
最后验证 \(\mathrm{Or}\) 对应 \(\mathrm{sgn}\)：\(\mathfrak S_k\) 在 \(\mathrm{Bij}([k],S)\) 上的预合成作用在商集 \(\mathrm{Bij}([k],S)/\mathfrak A_k\) 上恰给出一个两点轨道空间，其作用即由 \(\mathrm{sgn}\) 决定。证毕。
\end{proof}

\begin{corollary}[边界 uplift 的取向 parity：从二层 sheet parity 到三层二值覆盖]\label{cor:bdry-orientation-parity-uplift}
沿用 \(\Delta_m^{\mathrm{bdry}}\) 的 boundary uplift 位移集合记号（例如表 \ref{tab:foldbin_boundary_lift_m6_m8} 给出 \(m\in\{6,7,8\}\) 的证书输出），并对每个 boundary 纤维 \(w\) 令 \(\Delta_m^{\mathrm{bdry}}(w)\) 表示其 uplift-delta 作为无序有限集。
在“不引入外加寄存器且对分支重标号保持自然性”的约束下，可定义一个跨层数统一的二值跳跃对象：
\[
\Pi_m(w):=\mathrm{Or}\!\bigl(\Delta_m^{\mathrm{bdry}}(w)\bigr)\in \mathbf{Tors}(\ZZ_2).
\]
则：
\begin{enumerate}
\normalfont
  \item \(\Pi_m(w)\) 对 uplift 分支的任意重标号（置换）天然不变，并随双射函子性变化；
  \item 当 \(|\Delta_m^{\mathrm{bdry}}(w)|=2\) 时（例如 \(m=6\) 的 boundary 二层覆盖），因 \(\mathfrak A_2=\{e\}\) 有规范同构 \(\mathrm{Or}(S)\cong \mathrm{Bij}([2],S)\cong S\)，故 \(\Pi_m(w)\) 退化回二层 sheet parity 的等价表述；
  \item 当 \(|\Delta_m^{\mathrm{bdry}}(w)|=3\) 时（例如 \(m=7,8\) 的三层覆盖），\(\Pi_m(w)\) 给出唯一非平凡、且无需指定“特异层”的二值跳跃对象；其对象类型为取向 torsor，而非“层标签上的 \(\ZZ_2\)-分级”。
\end{enumerate}
\end{corollary}

\begin{proof}
第 (1) 点是 \(\mathrm{Or}\) 的函子性。
第 (2) 点由 \(\mathfrak A_2=\{e\}\) 直接推出。
第 (3) 点由定理 \ref{thm:bdry-binary-jump-orientation-functor-uniqueness} 在 \(k=3\) 的唯一性分类给出：置换自然且非平凡的二值结构在同构意义下只能是取向 torsor。证毕。
\end{proof}

\begin{proposition}[fold-gauge 的符号降维：二值跳跃从中心迁移到交换化]\label{prop:bdry-fold-gauge-sign-abelianization}
对任意有限集合 \(S\)（\(|S|=d\ge 2\)），令
\[
\mathrm{Lab}(S):=\mathrm{Bij}([d],S)
\]
为其标号集合，则 \(\mathrm{Lab}(S)\) 为主 \(\mathfrak S_d\)-torsor。
记 \(\mathrm{sgn}:\mathfrak S_d\to\ZZ_2\) 为符号同态，则存在规范同构
\[
\mathrm{Or}(S)\ \cong\ \mathrm{Lab}(S)\times_{\mathfrak S_d,\mathrm{sgn}}\ZZ_2.
\]
特别地：当 \(d=2\) 时 \(\mathfrak S_2\cong\ZZ_2\) 且其中心即自身，二值跳跃可实现为纤维内的“中心型”层翻转；而当 \(d\ge 3\) 时 \(Z(\mathfrak S_d)=\{e\}\) 但 \(\mathfrak S_d^{\mathrm{ab}}\cong\ZZ_2\)，故二值跳跃只能通过交换化（符号表示）出现，即表现为取向 torsor 而非层标签。
\end{proposition}

\begin{proof}
由定义 \(\mathrm{Or}(S)=\mathrm{Bij}([d],S)/\mathfrak A_d\)。
另一方面，纤维积 \(\mathrm{Lab}(S)\times_{\mathfrak S_d,\mathrm{sgn}}\ZZ_2\) 等价于在 \(\mathrm{Lab}(S)\) 中把相差一个偶置换的标号识别为同类，故与上述商完全一致。
关于中心与交换化的断言为对称群的基本结构：\(Z(\mathfrak S_d)=\{e\}\)（\(d\ge 3\)）且 \(\mathfrak S_d^{\mathrm{ab}}\cong\ZZ_2\)（\(d\ge 2\)）。证毕。
\end{proof}

\endinput
